% LexBridge — ACL-style skeleton. Compile with the ACL style files (acl.sty) placed next to this file.
\documentclass[11pt]{article}
\usepackage[review]{acl}
\usepackage{times,latexsym,booktabs,amsmath}
\usepackage[T1]{fontenc}
\usepackage[utf8]{inputenc}

\title{LexBridge: Code-Aware Cross-Code Retrieval for India's Criminal Law Transition}
\author{Anonymous submission}

\begin{document}
\maketitle

\begin{abstract}
On 1 July 2024 India replaced the Indian Penal Code (IPC), Code of Criminal Procedure and Evidence Act with
the Bharatiya Nyaya Sanhita (BNS), Bharatiya Nagarik Suraksha Sanhita and Bharatiya Sakshya Adhiniyam.
Questions are now asked in BNS terms while precedent cites the IPC, and identical section numbers denote
different offences across codes (IPC~302 is murder; BNS~302 is wounding religious feelings).
We present LexBridge, a retrieval-augmented QA system with (C1) a code-aware section resolver, (C2) cross-code
graph retrieval, (C3) precedent transferability scoring, (C4) a temporal applicability gate, and (C5)
LexBridge-QA, a benchmark targeting these failure modes. % TODO: headline numbers on the annotated benchmark.
\end{abstract}

\section{Introduction}
% Motivation: the transition; four failure modes of plain RAG (collision, cross-code gap, transferability, temporal).
% Contributions C1--C5.

\section{Related Work}
% Legal IR / RAG (COLIEE, AILA, ILDC/CJPE), statute retrieval, legal NER, temporal reasoning in law,
% legislative change and versioned statutes, hallucination and citation verification in legal LLMs.

\section{Problem Formulation}
% Codes, sections, mapping M: old -> new ∪ {⊥}; change types; applicability tags; governing code by offence date.

\section{LexBridge}
\subsection{C1: Code-Aware Section Resolution}
% NER formats; six-rule priority with confidences.
\subsection{C2: Cross-Code Graph Retrieval}
% Graph schema; hop; hybrid BM25+dense+graph with RRF; collision filter.
\subsection{C3: Precedent Transferability}
% Rule layer (diff ratio, clause split) + LLM classifier; label→tag map.
\subsection{C4: Temporal Gate}
% Offence-date extraction; governing family; dual-code answers when undated.
\subsection{Generation and Verification}
% Citation format, dual citation, claim verification and single regeneration.

\section{LexBridge-QA}
% Templates, candidate generation, annotation protocol, κ, statistics per category.

\section{Experimental Setup}
% Corpora (ILDC, Indian Kanoon sample), mappings (HF + IEEE, conflicts), models (bge-m3, bge-reranker, qwen2.5-7b),
% baselines (No-RAG, BM25, Dense, Hybrid), metrics (CER@k, code acc, R@k, MRR, nDCG@k, transfer F1, temporal acc,
% citation P/R, faithfulness).

\section{Results}
% NOTE: tables below are currently auto-generated from fixture data + a mock LLM (pipeline check, not results).
% AUTO-GENERATED by lexbridge.eval.ablations. Fixture corpus + deterministic mock LLM:
% these numbers check that the pipeline runs; they are NOT research results.

\begin{table*}[t]
\centering\small
\begin{tabular}{lcccccccccc}
\toprule
System & CER@5$\downarrow$ & Code acc & R@5 & MRR & nDCG@5 & Transfer F1 & Temporal acc & Cit.~P & Cit.~R & Faith. \\
\midrule
No-RAG & -- & 0.558 & 0.000 & 0.000 & 0.000 & 0.194 & 0.286 & 0.000 & 0.000 & -- \\
BM25-RAG & 1.000 & 1.000 & 0.729 & 0.630 & 0.587 & 0.194 & 0.286 & 0.343 & 0.679 & 1.000 \\
Dense-RAG & 0.923 & 0.977 & 0.695 & 0.631 & 0.577 & 0.194 & 0.214 & 0.329 & 0.633 & 1.000 \\
Hybrid-RAG & 1.000 & 1.000 & 0.714 & 0.638 & 0.578 & 0.194 & 0.286 & 0.343 & 0.671 & 1.000 \\
\textbf{LexBridge} & 0.000 & 1.000 & 0.957 & 0.986 & 0.961 & 0.931 & 1.000 & 0.881 & 0.867 & 1.000 \\
\bottomrule
\end{tabular}
\caption{Retrieval results on LexBridge-QA ($n=50$, fixture corpus, mock LLM; not research results).}
\label{tab:main}
\end{table*}

% AUTO-GENERATED by lexbridge.eval.ablations. Fixture corpus + deterministic mock LLM:
% these numbers check that the pipeline runs; they are NOT research results.

\begin{table*}[t]
\centering\small
\begin{tabular}{lcccccccccc}
\toprule
System & CER@5$\downarrow$ & Code acc & R@5 & MRR & nDCG@5 & Transfer F1 & Temporal acc & Cit.~P & Cit.~R & Faith. \\
\midrule
full & 0.000 & 1.000 & 0.957 & 0.986 & 0.961 & 0.942 & 1.000 & 0.881 & 0.867 & 1.000 \\
$-$C1 resolver & 0.154 & 0.977 & 0.924 & 0.917 & 0.889 & 0.942 & 1.000 & 0.700 & 0.767 & 1.000 \\
$-$C2 graph hop & 0.154 & 1.000 & 0.814 & 0.810 & 0.751 & 0.942 & 1.000 & 0.595 & 0.710 & 1.000 \\
$-$C3 transferability & 0.000 & 1.000 & 0.957 & 0.986 & 0.961 & 0.176 & 1.000 & 0.881 & 0.867 & 1.000 \\
$-$C4 temporal gate & 0.077 & 0.977 & 0.948 & 0.971 & 0.947 & 0.942 & 0.500 & 0.852 & 0.838 & 1.000 \\
$-$role$-$aware chunking & 0.000 & 1.000 & 1.000 & 0.986 & 0.989 & 0.942 & 1.000 & 0.564 & 1.000 & 1.000 \\
$-$verifier & 0.000 & 1.000 & 0.957 & 0.986 & 0.961 & 0.942 & 1.000 & 0.881 & 0.867 & 1.000 \\
\bottomrule
\end{tabular}
\caption{Ablations of LexBridge components ($n=50$, fixture corpus, mock LLM).}
\label{tab:ablation}
\end{table*}


\section{Error Analysis}
% Mapping errors, resolver failures on unlabelled numbers, transferability disagreements (e.g. IPC 302 → BNS 103(1)/(2)).

\section{Limitations}
Criminal law only; mapping errors propagate into retrieval and applicability tags; rule-based role tagging;
procedural savings under the BNSS are flagged but not modelled. LexBridge does not provide legal advice.

\section*{Ethics Statement}
% Risks of over-reliance by litigants; data licensing (India Code, ILDC terms, Indian Kanoon terms); annotator compensation.

\bibliography{custom}
\end{document}
